\documentclass[10pt,a4paper]{article}

\usepackage[margin=1.55cm]{geometry}
\usepackage{fontspec}
\usepackage{xcolor}
\usepackage{hyperref}

\setmainfont{Times New Roman}

\definecolor{linkblue}{RGB}{20,75,140}
\hypersetup{
  colorlinks=true,
  linkcolor=black,
  urlcolor=linkblue
}

\setlength{\parindent}{0pt}
\setlength{\parskip}{2pt}
\linespread{1.02}
\setlength{\leftmargini}{1.4em}
\let\originalitemize\itemize
\let\endoriginalitemize\enditemize
\renewenvironment{itemize}
  {\originalitemize
   \setlength{\itemsep}{1.8pt}
   \setlength{\topsep}{2.5pt}
   \setlength{\parsep}{0pt}
   \setlength{\partopsep}{0pt}}
  {\endoriginalitemize}

\newcommand{\cvname}[1]{{\LARGE\bfseries #1}\par}
\newcommand{\cvsection}[1]{
  \vspace{0.55em}{\large\bfseries #1}\par
  \vspace{0.1em}\hrule\vspace{0.35em}
}
\newcommand{\entry}[4]{
  \textbf{#1}\hfill #2\\
  \textit{#3}\hfill\textit{#4}
}
\newcommand{\project}[2]{
  \textbf{#1}\\[-1pt]
  \textit{#2}
}

\pagestyle{empty}

\begin{document}

\cvname{Bangzhe Huang}
\vspace{0.12em}
\href{mailto:huangbzphys@gmail.com}{huangbzphys@gmail.com}
\quad|\quad
\href{https://github.com/bangzhehuang}{github.com/bangzhehuang}
\quad|\quad
\href{https://bangzhehuang.github.io/}{bangzhehuang.github.io}

\vspace{0.4em}
\textbf{Research interests:} automated research, research world models,
AI4AI, scientific machine learning, and many-body physics

\cvsection{Education}

\entry
  {Fudan University, Department of Physics}
  {Shanghai, China}
  {B.Sc. in Physics, incoming undergraduate}
  {September 2026 -- expected June 2030}

\cvsection{Research Experience}

\entry
  {College of AI, Tsinghua University}
  {Beijing, China}
  {Research Intern, Prof. Ziming Liu's group}
  {Current}

\vspace{0.15em}
\project
  {AI4AI: Learning to Design a Meta-model}
  {Independent research; active pilot}

\begin{itemize}
  \item Built a controlled AI4AI pilot in which one meta-model predicts
  training dynamics and a second model learns architecture-dependent test
  log-MAE to guide the design of the first model.
  \item Defined an exhaustive DAG search space from four homogeneous MLP blocks
  and one softmax gate, generated 10,006 canonical candidates, and built a graph
  encoder for supervised architecture-to-outcome prediction.
  \item Implemented the candidate-training and meta-model pipelines and an
  interactive interface for manual DAG construction, outcome prediction, and
  automatic architecture design. Preliminary experiments surface a
  ResNet-like design.
\end{itemize}

\project
  {Physics-Informed Many-Body Foundation Models: Scalable Optimization across
  Heterogeneous Tasks}
  {Project lead; advisor: Prof. Jing Wang, Fudan University; preprint in preparation}

\begin{itemize}
  \item Proposed a task-modular many-body foundation model that preserves each
  Hamiltonian family's physical priors and ansatzes while learning
  high-level representations for cross-task composition.
  \item Led problem formulation, PyTorch implementation, and experimental
  evaluation against unified joint-training baselines across task scales,
  task combinations, and fine-tuning settings.
  \item Preliminary results show comparable generalization and adaptation when
  unified training is optimizable, while the modular system remains trainable
  in regimes where the unified baseline becomes ill-conditioned.
\end{itemize}

\project
  {Open-Ended Scientific Research Benchmark for Physics Research Agents}
  {Independent research; academic discussion with Prof. Siheng Chen, Shanghai Jiao Tong University}

\begin{itemize}
  \item Designed open-ended physics research tasks with stable ground truth
  defined by low-dimensional physical structure and equivalent observables,
  without prescribing a single valid research trajectory.
  \item Built an end-to-end pilot covering task interpretation, simulation
  design, code generation and execution, feedback-driven action selection, and
  final physical auditing.
\end{itemize}

\project
  {Phase Transport and Defect Dynamics in Nonreciprocal Spin Systems}
  {Lead researcher; advisor: Prof. Zi Cai, Shanghai Jiao Tong University}

\begin{itemize}
  \item Used Monte Carlo simulations to study the microscopic origin of
  macroscopic oscillations in two-dimensional nonreciprocal spin systems and
  extracted oscillation frequency, domain-wall density, and defect observables.
  \item Proposed the effective transport relation
  $\omega=\rho v_{\mathrm{eff}}$ and identified two dynamical regimes associated
  with periodic wall expansion and vortex-constrained topological structures.
  \item Tested the resulting physical picture across lattice structures and
  symmetric bond disorder.
\end{itemize}

\cvsection{Technical Skills}

\textbf{Machine learning:} PyTorch, deep learning, scientific foundation
models, graph neural networks, LLM agents, benchmark design\\
\textbf{Scientific computing:} Python, Monte Carlo simulation, numerical
experiments, statistical analysis\\
\textbf{Physics:} quantum mechanics, statistical physics, many-body physics

\end{document}
